% !TeX root = ../main.tex
\section{Background and Related Work}
\label{sec:background-related-work}

\subsection{Background}
\label{sec:background}

\emph{Architectural Design Decisions and ADRs.}
Software architecture is widely characterized as the composition of fundamental
design decisions that shape a system's structure, behavior, and quality
attributes~\cite{jansen2005decisions,kruchten2006architectural}.
Capturing the context, trade-offs, and rationale behind these choices is critical
for system evolution, knowledge transfer, and preventing architectural drift.
Architectural Decision Records (ADRs) provide a lightweight, structured
documentation format to record individual decisions~\cite{nygard2011adr}.
In Nygard's established formulation, an ADR captures a title, status (e.g.,
Proposed, Accepted, Superseded), business and technical context, the decision
made, and its resulting consequences (both positive and negative trade-offs).
Despite their recognized value in maintaining architectural traceability,
consistent ADR adoption in industrial practice is frequently hampered by time
constraints, high cognitive overhead, and the burden of manual authoring.

\emph{Architectural Patterns for AI-Enabled Systems.}
Architectural patterns offer reusable, proven solutions to recurring design
problems within specific operational contexts. While traditional software
patterns address general modularity, loose coupling, and component
interaction, AI-enabled systems introduce distinct engineering concerns.
These include managing non-deterministic model inference, training and
serving pipeline orchestration, data drift, continuous monitoring, and
specialized compute infrastructure.
In response, recent software engineering research has codified specialized
catalogs and taxonomies of design patterns tailored for machine learning and
AI-based applications~\cite{washizaki2019patterns,heiland2023patterns}.
These patterns span key lifecycle stages, including data acquisition,
feature preparation, model training, model deployment, serving, and runtime
observability~\cite{aiPatternsRepository}.
However, matching concrete system requirements to appropriate AI patterns
remains an unstructured, manual task that requires deep familiarity with both
the operational domain and extensive pattern repositories.

\emph{LLMs in Architecture: Retrieval versus Generative Reasoning.}
The emergence of Large Language Models (LLMs) with advanced natural language
understanding and reasoning capabilities presents new opportunities for
architectural decision support. In this context, two primary paradigms exist:
\begin{enumerate}
    \item \emph{Dense Semantic Retrieval (RAG):} Retrieval-Augmented
    Generation techniques embed requirement statements and pattern definitions
    into continuous vector spaces to retrieve top-ranking candidates based on
    semantic proximity~\cite{lewis2020rag}. Dense retrieval is computationally
    efficient and scalable across large catalogs, but operates primarily on
    lexical and semantic surface similarity without evaluating domain-specific
    trade-offs, structural constraints, or verifiable requirement satisfaction.
    \item \emph{Generative Applicability and Decision Synthesis:} Prompting
    frontier LLMs with explicit requirement contexts, acceptance criteria,
    and structured pattern specifications enables direct reasoning over
    compatibility. Beyond simple similarity, generative models can categorize
    applicability under formalized decision rules, cite explicit textual
    evidence from requirements, and synthesize complete, structured ADR drafts
    documenting context and trade-offs.
\end{enumerate}
Understanding the operational relationship, overlap, and practical boundaries
between dense semantic retrieval and generative reasoning is central to designing
effective, requirement-grounded architectural assistance pipelines.

\subsection{LLM-Assisted Software Architecture Design}
\label{sec:related-work-architecture}

Recent research has increasingly investigated the capabilities and boundaries
of LLMs as intelligent assistants in software architecture design.
Soliman and Keim~\cite{soliman2025architectural} conducted an exploratory
empirical evaluation probing whether large language models embody reliable
software architectural knowledge (AK), assessing GPT responses against
established ground truth on the large-scale HDFS system. While the model
demonstrated moderate recall in retrieving general architectural concepts, it
exhibited substantially lower precision, particularly when tasked with identifying
quality attribute solutions and capturing nuanced design rationales.
Their findings established that while LLMs can serve as an exploratory source
for initial architectural ideas, independent expert verification remains
indispensable due to the absence of reliable semantic self-grounding.

Beyond unstructured knowledge queries, researchers have explored structured,
methodology-guided interactions between human architects and generative models.
Cervantes et al.~\cite{cervantes2026assistants} examined the broader framing
of LLMs as collaborative architectural assistants, identifying their strengths
in brainstorming design concepts, structuring documentation, and exploring
technology trade-offs. To operationalize this collaboration within a formal
architectural method, Cervantes et al.~\cite{cervantes2025add} proposed an
LLM-assisted approach embedded within the Attribute-Driven Design (ADD) framework.
By supplying the model with an explicit ADD method description, an architect persona,
and a structured iteration plan, the system guides the LLM to co-create
architectural design artifacts. Their case-study evaluations with professional
architects revealed that while generated designs partially satisfy architectural
drivers, continuous human oversight, explicit decision boundaries, and iterative
refinement are essential to prevent ungrounded architectural drift.

These studies demonstrate the feasibility of employing LLMs across diverse
architectural tasks, including conceptual design, quality-attribute reasoning,
and structured method execution. However, existing work has predominantly focused
on open-ended conversational prompting, broad architectural drivers, or manual
prompt-response loops across general software systems.
Systematic, requirement-grounded pattern selection across codified catalogs of
AI-specific architectural patterns---coupled with the automated synthesis of
traceable decision records---remains underexplored.
Our work addresses this gap by formalizing a requirement-by-pattern census
evaluation that enforces strict textual grounding rules, compares generative
reasoning against dense semantic retrieval baselines, and generates verifiable,
structured ADRs for professional evaluation.

\subsection{LLM-Based ADR and Design Rationale Generation}
\label{sec:related-work-adr}

Documenting architecture decisions through structured records (ADRs) and
capturing the underlying design rationale (DR) are central challenges in
architectural knowledge management. Because manual documentation is often
neglected in fast-paced software development, recent studies have investigated
whether LLMs can automate or support decision authoring.

Dhar et al.~\cite{dhar2024generating} conducted an exploratory empirical study
evaluating the feasibility of generating architectural decisions using LLMs.
Focusing on existing open-source ADR corpora, they formulated the task as
generating the \emph{Decision} section of an ADR conditioned on a human-authored
\emph{Context} section. They evaluated zero-shot, few-shot, and fine-tuning
approaches across models including GPT-4, GPT-3.5, and Flan-T5, using automated
lexical and semantic similarity metrics (BLEU, ROUGE, BERTScore) against
human-written reference decisions. Their results showed that advanced models like
GPT-4 can generate plausible decisions in zero-shot settings, but highlighted
key limitations: lexical overlap metrics correlate poorly with architectural
soundness, and the generation pipeline presupposes that an architect has
already synthesized and documented an explicit decision context.

Focusing on the argumentative reasoning behind architectural choices,
Zhou et al.~\cite{zhou2026rationale} investigated LLM generation of design
rationale across 100 architecture-related problems mined from Stack Overflow
and GitHub issues and discussions. Evaluating five LLMs across zero-shot,
Chain-of-Thought (CoT), and LLM-based agent strategies against human-expert
ground truth, they observed moderate recall (0.63--0.72) but low precision
(0.27--0.28) in generating trade-off arguments. While a substantial majority of
generated arguments not originally mentioned by experts were deemed helpful,
between 1.6\% and 3.2\% were potentially misleading or incorrect, underscoring
the risk of hallucinated trade-offs in automated rationale recovery.

These investigations provide important baseline evidence regarding LLM
competence in decision and rationale formulation. However, important differences
distinguish prior efforts from the workflow evaluated in this paper:
\begin{itemize}
    \item \emph{Input Grounding and Granularity:} Prior studies condition
    generation on pre-formulated ADR context paragraphs~\cite{dhar2024generating}
    or informal forum discussions~\cite{zhou2026rationale}. In contrast, our
    pipeline operates directly on structured industrial requirements (user
    stories and acceptance criteria) and evaluates applicability across an
    established catalog of codified architectural patterns.
    \item \emph{Artifact Completeness:} Rather than predicting an isolated
    field (such as the decision statement or rationale arguments), our
    generator produces complete, multi-field Nygard ADRs (Title, Status,
    Context, Decision, and Consequences), maintaining internal consistency
    between cited requirement evidence and pattern trade-offs.
    \item \emph{Pipeline Traceability:} We evaluate ADR generation as an
    explicit downstream stage conditioned on a verified applicability
    assessment, establishing an unbroken audit trail from functional
    requirements to candidate pattern selection and final decision records.
\end{itemize}

\subsection{Positioning Against AI-System Pattern Research}
\label{sec:related-work-positioning}

The codification and operational use of architectural design patterns for
AI-enabled systems represent an active and rapidly expanding research frontier.
Closest to our investigation, Oliveira et al.~\cite{oliveira2026designpartners}
evaluated the capacity of state-of-the-art LLMs (Gemini 1.5 Flash, Claude 3.5 Sonnet,
and o1-mini) to serve as ``design partners'' by generating architectural pattern
solutions across 65 documented AI patterns from the SWE4AI repository~\cite{aiPatternsRepository}.
Using a mixed-methods evaluation combining cosine similarity, expert review of
divergent solutions, and blind practitioner preference, they observed that LLMs
can synthesize conceptual pattern solutions that match or surpass literature-defined
descriptions in over 90\% of cases.

While both studies investigate architectural patterns for AI-enabled systems using
overlapping pattern taxonomies derived from the SWE4AI corpus~\cite{heiland2023patterns,aiPatternsRepository},
their research objectives, inputs, and evaluation tasks diverge fundamentally:
\begin{enumerate}
    \item \emph{Pattern Generation versus Requirement-Grounded Selection:}
    Oliveira et al. evaluated \emph{solution synthesis in the abstract}: given a
    pattern's documented problem statement and intent, can an LLM generate an
    equivalent or improved pattern structure? In contrast, ADR4AI evaluates
    \emph{requirement-driven applicability}: given concrete, uncurated industrial
    user stories and acceptance criteria from four production systems, can a model
    discern which existing, codified patterns actually apply under strict textual
    evidence rules?
    \item \emph{Architectural Artifact Production:} Whereas Oliveira et al. focused
    on conceptual pattern descriptions, our workflow connects pattern selection
    directly to the automated authoring of complete Architectural Decision Records
    (ADRs) documenting operational context, chosen solutions, and positive and
    negative trade-offs.
    \item \emph{Retrieval versus Reasoning Comparison:} While prior pattern
    generation research assessed LLM generative capacity in isolation, our study
    empirically compares generative applicability reasoning against a standard dense
    semantic retrieval (RAG) baseline, quantifying overlap and divergence across
    an exhaustive census of 840 story-pattern pairs.
\end{enumerate}

From a complementary perspective, automated pattern identification has also been
investigated in retrospective code analysis. J{\'a}nki and Bilicki~\cite{janki2023recognition}
demonstrated that GPT models guided by formal rule systems can achieve 90\%
accuracy in detecting Model-View-Whatever (MVW) architectural patterns within
Angular web application source code. Their work addresses pattern \emph{recognition}
in existing software artifacts based on deterministic structural rules. By contrast,
ADR4AI addresses \emph{prospective selection and architectural decision documentation}
during requirements engineering and early design, before code artifacts exist.

\emph{Research Gap.}
In summary, while prior research has probed LLM architectural knowledge~\cite{soliman2025architectural},
method-guided design assistance~\cite{cervantes2026assistants,cervantes2025add},
isolated ADR decision completion~\cite{dhar2024generating}, rationale argument
recovery~\cite{zhou2026rationale}, abstract pattern solution synthesis~\cite{oliveira2026designpartners},
and code-level pattern recognition~\cite{janki2023recognition}, no prior study provides:
(i) an empirical evaluation of requirement-grounded architectural pattern selection
over a codified catalog of AI-enabled design patterns;
(ii) a systematic comparative analysis between generative LLM reasoning and dense
semantic retrieval across an exhaustive census of industrial requirements; and
(iii) an end-to-end, traceable pipeline that transforms requirement-grounded pattern
selections into complete, inspectable Architectural Decision Records evaluated by
practicing software architects.
This study directly addresses this research gap.



